\chapter{Conclusão e Trabalho Futuro}
\label{cp:conclusao}

O desenvolvimento do \textit{Plot in a Pot} (\gls{piap}) teve como premissa a resolução de um problema real e recorrente enfrentado por autores de \gls{ni} e \textit{game designers}: a complexidade geométrica e semântica que surge durante o planeamento de histórias não-lineares, agravada pela falta de ferramentas integradas para gestão de localização multi-idioma e auditoria estática de fluxo.

\section{Conclusão}
\label{sec:conclusao-recap}

Com a conclusão deste projeto informático, todos os objetivos gerais e específicos inicialmente traçados no Capítulo~\ref{cp:introduction} foram integralmente alcançados:
\begin{itemize}
    \item \textbf{Interface e Acessibilidade Visual:} Foi desenvolvido um editor gráfico inovador assente numa estética neo-brutalista de alto contraste. Esta opção de \textit{design} provou ser não apenas esteticamente diferenciadora, mas de extrema pertinência funcional para utilizadores com limitações de visão (p. ex. daltónicos), conforme validado pelos testes com o participante U6.
    \item \textbf{Gestão de Localização:} A implementação da matriz de tradução bidimensional baseada em grelhas dinâmicas, com importação e exportação em formato \gls{csv} em conformidade com a norma RFC~4180, eliminou por completo a necessidade de gerir múltiplos ficheiros de texto externos ou injetar strings rígidas (\textit{hardcoded}) nas passagens do Twine.
    \item \textbf{Validação e Algoritmia:} A adoção de uma estrutura de dados assente em listas de adjacências bidirecionais suportou com sucesso a execução em tempo real de análises estáticas baseadas em explorações via \gls{bfs}. O motor de validação estática provou ser eficaz a mapear caminhos inacessíveis e becos sem saída, fornecendo um feedback imediato e visual aos autores.
    \item \textbf{Simulação Automatizada:} O modo de simulação (\textit{PlayMode}) foi complementado pelo \textit{Smart Walker}, um agente inteligente baseado em procuras por novidade (\textit{Novelty Search}) capaz de realizar o percurso autónomo e teste lógico de narrativas complexas sem intervenção humana.
    \item \textbf{Suavização Geométrica:} O desenvolvimento das equações paramétricas das Splines Cúbicas de Catmull-Rom garantiu arestas interativas suaves com \textit{waypoints} para organização e otimização espacial do canvas de grafos.
\end{itemize}

Os testes estruturados com utilizadores de diferentes perfis revelaram taxas de sucesso muito elevadas na realização de tarefas (com uma média global de 95,8\%), demonstrando que a barreira de abstração erguida pela nossa plataforma é intuitiva quer para perfis puramente artísticos (autores/escritores), quer para utilizadores com forte background técnico (programadores).

\subsection{Porquê Escolher o Plot in a Pot?}
\label{subsec:porque-escolher-piap}
Com base nas limitações identificadas nas ferramentas de autoria existentes (analisadas na Secção~\ref{sec:analise-comparativa}), o \gls{piap} destaca-se como uma solução unificada que resolve falhas críticas de usabilidade e integridade no desenvolvimento de narrativas interativas:
\begin{itemize}
    \item \textbf{Fusão de Paradigmas (Design Visual com Validação Lógica):} Ao contrário do Twine (que carece de testes lógicos nativos e robustos) e do ChoiceScript/Ink (que operam puramente sob interfaces textuais), o \gls{piap} une a facilidade de arrastar e interligar nós visuais com a solidez de um validador estático e simulador inteligente automatizado.
    \item \textbf{Localização Nativa sem Dispersão de Dados:} Resolve a dispersão de ficheiros externos e a escrita rígida (\textit{hardcoded}) de diálogos através de uma matriz de localização baseada em grelhas dinâmicas, acelerando o fluxo de internacionalização com importação/exportação CSV.
    \item \textbf{Acessibilidade Inclusiva Funcional:} A interface neo-brutalista de alto contraste e a redundância de traços nas arestas mitigam barreiras cognitivas e visuais severas, permitindo que autores daltónicos desenhem e validem histórias de forma totalmente autónoma.
    \item \textbf{Soberania de Dados e Interoperabilidade:} Ao adotar a especificação aberta Twee 3 (SugarCube), o autor mantém controlo absoluto e livre sobre a sua obra, garantindo compatibilidade bidirecional síncrona com outros compiladores oficiais (como o Tweego).
\end{itemize}

\section{Trabalho Futuro}
\label{sec:trabalho-futuro}

Apesar dos excelentes resultados e estabilidade demonstrados, o \textit{Plot in a Pot} apresenta oportunidades claras de evolução, destacando-se a IA local e a edição colaborativa como prioridades imediatas:

\begin{enumerate}
    \item \textbf{[Prioridade Imediata] Execução de IA Local via \gls{webgpu}:} Substituir a dependência de APIs externas ou do Ollama local pela execução de \glspl{llm} leves diretamente no navegador (ex: WebLLM). Isto elimina a necessidade de instalações adicionais pelo utilizador, garantindo a soberania absoluta dos dados (RNF04).
    \item \textbf{[Prioridade Imediata] Edição Colaborativa em Tempo Real:} Integrar algoritmos de sincronização baseados em \glspl{crdt} (como Yjs ou Automerge) para permitir que múltiplos autores editem o grafo narrativo em simultâneo, mimetizando a experiência de ferramentas como o Figma.
    \item \textbf{Empacotamento Nativo para Desktop:} Exportar a aplicação para executáveis nativos (Windows, macOS e Linux) usando Tauri ou Electron, facilitando uma integração mais profunda com o sistema de ficheiros local.
    \item \textbf{Análise Formal com Model Checking:} Converter o grafo e as condições SugarCube em especificações formais (ex: Promela ou UPPAAL). Isto permitirá usar \textit{Model Checking} para provar matematicamente propriedades de segurança e vivacidade, como a garantia de caminhos viáveis até aos desfechos da história.
    \item \textbf{Módulo de Gestão de Personagens (\textit{Character Tracker}):} Adicionar um sistema integrado para mapear a presença e estados emocionais das personagens diretamente nos nós do canvas (React Flow), respondendo a uma sugestão recolhida nas sessões de avaliação.
    \item \textbf{Filtragem de Nós Privados na Exportação:} Implementar um filtro ativo que impeça a inclusão de passagens marcadas como privadas ou em desenvolvimento no ficheiro final compilado (Twee/HTML), prevenindo \textit{spoilers} ou a fuga de conteúdos inacabados.
\end{enumerate}

Em suma, o desenvolvimento do \textit{Plot in a Pot} alcançou com sucesso os objetivos propostos, demonstrando que a aproximação entre a engenharia de software, a teoria de grafos e a escrita criativa é capaz de simplificar e enriquecer a autoria de ficção interativa. Ao aliar a robustez da validação formal de espaço de estados à flexibilidade da interação visual e à acessibilidade inclusiva, a plataforma assume-se como uma ferramenta sólida para a modelação de narrativas não-lineares, servindo de ponto de partida e base de trabalho para investigações e desenvolvimentos futuros no domínio dos jogos interativos.
